\section{Methodology}
\subsection{Part 1: Hopper Continuous Control}

Part 1 focused on learning a locomotion policy for the Hopper environment using policy-gradient methods.
Three approaches were implemented and compared: REINFORCE, REINFORCE with baseline, and Actor-Critic architectures \cite{SuttonBarto2018}.

REINFORCE \cite{Williams1992} is the simplest policy algorithm, updating the policy parameters at the end of each episode using the full Monte Carlo return. The loss minimized at each update is:

\[
L = -\sum_t \left(G_t - b\right) \log \pi_\theta(a_t | s_t)
\]



where $G_t$ is the discounted return from timestep $t$, $b$ is a baseline and $G_t = \sum_{k=0}^{T-t} \gamma^k r_{t+k}$. This formulation ensures that the gradient estimator is unbiased for any choice of $b$, because $\mathbb{E}[\nabla_\theta \log \pi_\theta \cdot b] = 0$ for any constant. However, it has high variance due to the stochasticity of full-episode returns. In the first variant $b=0$ we have vanilla REINFORCE. In the second, the return is shifted by a constant baseline $b = 20$.
Both variants share the same policy network and training configuration: a two-layer MLP with hidden sizes [64, 64] and Tanh activations, trained using the Adam optimizer with learning rate $3 \times 10^{-4}$ and discount factor $\gamma = 0.99$.

We used orthogonal initialization \cite{saxe2013exact} instead of the standard normal weight initialization for better training stability. The hidden layers used a gain of $\sqrt{2}$ to preserve the gradient norms and avoid exploding or vanishing gradients. Furthermore, the output layer of the actor was initialized with a low gain of 0.01. This makes sure the initial policy outputs action values close to zero, stopping extreme movements that would destabilize the Hopper right away. Therefore, the early exploration is entirely governed by the learned standard deviation, which allows the agent to survive for a longer period of time and collect better trajectories during the initial training stage \cite{andrychowicz2020matters}.

The final implementation uses an Actor-Critic architecture, where the actor learns the policy while the critic estimates
the value function. The actor updates the policy based on the estimated advantages, modifying the probability of actions
based on how they turned out to be with respect to the critic’s expectations:

\[
L_{\text{actor}} = -\mathbb{E}\left[\hat A_t \log \pi_\theta(a_t|s_t)\right].
\]

In parallel, the critic minimizes the squared difference between the predicted state value and the temporal-difference target:

\[
L_{\text{critic}} = \mathbb{E}\left[\left(r_t+\gamma V_\phi(s_{t+1})-V_\phi(s_t)\right)^2\right].
\]

Since we are considering a continuous environment, all models' policy outputs both the mean and standard
deviation of a Gaussian action distribution:

\[
a_t \sim \mathcal N(\mu_\theta(s_t), \sigma_\theta(s_t)),
\]

where both parameters are learned by the policy network. Actions are sampled during training to encourage exploration,
while the mean action $\mu_\theta(s_t)$ is used during evaluation.

To improve the quality of the policy updates, Generalized Advantage Estimation (GAE)
was implemented \cite{SchulmanGAE2015}, with advantages computed as:

\[
\hat A_t^{GAE} = \sum_{l=0}^{\infty} (\gamma\lambda)^l \delta_{t+l}
\]

where

\[
\delta_t = r_t + \gamma V(s_{t+1}) - V(s_t)
\]

is the temporal-difference error.


Subsequently, advantages were normalized before policy updates in order to stabilize gradients,
improve critic's estimates and reduce training instability. A minimum standard deviation (sigma floor)
was enforced to prevent premature collapse of exploration and promote it even in later stages.
Additionally, a triggered learning-rate decay mechanism was introduced during later stages of training to reduce
update magnitude after the policy had discovered effective locomotion behaviors.


More specifically, the decay was activated only once one of the trigger conditions reported in Table~\ref{tab:lr_decay}
was satisfied, causing the learning rate to linearly reduce from $3 \times 10^{-4}$ to $10^{-5}$ until the end of training.

\begin{table}[ht]
\centering
\caption{Learning Rate Decay Triggers}
\label{tab:lr_decay}
\resizebox{\columnwidth}{!}{
\begin{tabular}{ll}
\hline
Trigger condition & Threshold \\
\hline
Rolling avg. reward \& episode length & avg\_reward\_100$\geq$1000, avg\_length\_100 $\geq$ 500 \\
Best rolling avg. reward & best\_avg\_reward $\geq$ 1500 \\
Episode fallback trigger & episode $\geq$ 45,000 \\
Initial learning rate & $3 \times 10^{-4}$ \\
Minimum learning rate & $10^{-5}$ \\
\hline
\end{tabular}
}
\end{table}

The goal was to reduce update magnitude once the Hopper had likely discovered jumping, as this behaviour leads to
higher rewards while being much more unstable. All thresholds were chosen empirically.

% high average reward together
% with long episode length meant sustained hopping was disovered, and also best\_avg\_reward $\geq$ 1500 indicated clear
% high-performance runs. The 45,000-episode condition acted as a fallback that would have forced the LR decay initialisation,
% again this threshold value was set based on previous runs, which showed that hopping usually emerged around 40k--45k episodes.

The final Actor-Critic configuration used the hyperparameters reported in Table~\ref{tab:ac_configuration}.
All values were selected empirically from previous preliminary training runs.


\begin{table}[ht]
\centering
\caption{Final Actor-Critic Configuration}
\label{tab:ac_configuration}
\resizebox{\columnwidth}{!}{
\begin{tabular}{ll}
\hline
Hyperparameter & Value \\
\hline
Discount Factor ($\gamma$) & 0.99 \\
GAE Lambda ($\lambda$) & 0.95 \\
Learning Rate & $3 \times 10^{-4}$ \\
Minimum Learning Rate After Decay & $1 \times 10^{-5}$ \\
Sigma Floor & 0.1 \\
Initial Learned Sigma & 0.5 (through Softplus + Sigma Floor) \\
Entropy Coefficient & 0.0 \\
Training Episodes & 70,000 \\
Learning Rate Decay Trigger & Performance-triggered with $\sim$45k episodes fallback \\
Hidden Layers & 64, 64 \\
Activation Function & Tanh \\
Optimizer & Adam \\
\hline
\end{tabular}
}
\end{table}



Performance was evaluated using episodic reward, rolling average reward (100-episode window),
episode length, policy entropy, learned action standard deviation, actor loss and critic loss.
























\subsection{Part 2: PandaPush, PPO, SAC and Domain Randomization}

Part 2 investigates sim-to-real transfer using the PandaPush-v3 robotic manipulation environment. The task consists of controlling a Franka Panda robotic arm to push a cube toward a desired target position.
To move the arm, the agent commands 3D Cartesian displacements of the end-effector (the seven joint angles are resolved by inverse kinematics), while the reward depends on the cube--goal distance.
Therefore, the agent must learn to position the end-effector so as to push the cube toward the target and achieve task success.
The observation space contains the robot state, cube state, and target position information. Actions are continuous Cartesian end-effector displacements, defined as:
\[
a_t=[dx,dy,dz].
\]

The dense reward function is the euclidean distance between the cube and target positions, defined as
\[
r_t = -\left\| p_{obj}(t)-p_{goal} \right\|_2,
\]
where $p_{obj}$ and $p_{goal}$ denote the cube and target positions respectively. A task is considered successful when the final cube-target distance is below 5 cm.

Two reinforcement-learning algorithms were evaluated: Proximal Policy Optimization (PPO) and Soft Actor-Critic (SAC).
Proximal Policy Optimization (PPO) is an on-policy actor-critic algorithm that limits policy updates through a clipped objective,
\[
L^{CLIP} = \mathbb E \Big[ \min( \rho_t \hat A_t, \text{clip}(\rho_t,1-\epsilon,1+\epsilon)\hat A_t ) \Big],
\]

where $\hat A_t$ is the advantage estimate, $\epsilon$ is a hyperparameter that controls the clipping range, and

\[
\rho_t = \frac{\pi_\theta(a_t|s_t)}{\pi_{\theta_{old}}(a_t|s_t)}
\]

is the probability ratio between the new and old policies.


The idea behind this algorithm is that excessively large policy updates result in unstable training, so stability is achieved to keep policy updates within a trust region (PPO is in fact the direct evolution of the Trust Region Policy Optimization (TRPO) \cite{SchulmanTRPO2015} algorithm).
The clipping mechanism improves stability but reduces sample efficiency because collected trajectories cannot be reused extensively.


Soft Actor-Critic (SAC) is an off-policy actor-critic method that maximizes both reward and policy entropy,
\[
J = \mathbb E[R] + \alpha H(\pi).
\]
SAC combines a replay buffer, twin critics to reduce value overestimation, and automatic entropy tuning. These characteristics generally provide better exploration and sample efficiency than PPO, especially in continuous-control robotic tasks.
Concretely, the twin critics minimise the soft Bellman residual
\[
J_Q(\theta)=\mathbb{E}\big[(Q_\theta(s_t,a_t)-(r_t+\gamma\,\mathbb{E}_{s_{t+1}}[V_{\bar\theta}(s_{t+1})]))^2\big],
\]
with $V(s)=\mathbb{E}_{a\sim\pi}[\min_i Q_i(s,a)-\alpha\log\pi(a|s)]$; the actor minimises
$J_\pi(\phi)=\mathbb{E}[\alpha\log\pi_\phi(a|s)-Q_\theta(s,a)]$, and the temperature $\alpha$ is
tuned automatically by minimising $J(\alpha)=\mathbb{E}[-\alpha(\log\pi(a|s)+\bar{\mathcal{H}})]$
against a target entropy $\bar{\mathcal{H}}=-\dim(\mathcal{A})$.

To improve transfer across different dynamics, Domain Randomization (DR) was applied by varying the cube mass during training.

Three configurations were evaluated:
\begin{itemize}
    \item No DR: fixed source mass $m=1$ kg.
    \item Uniform Domain Randomization (UDR): $m \sim U(0.5,8.0)$ kg, resampled at every episode reset; the range brackets both the source nominal ($1$ kg) and the target ($5$ kg) without being centered on the (assumed unknown) target mass.
    \item Automatic Domain Randomization (ADR): an AutoDR-style curriculum over the support $[0.5,8.0]$ kg. It starts from a narrow interval $\phi=[1.0-h,\,1.0+h]$, with $h=\min(0.1,\,0.25\,\Delta_{\text{range}})$,
    around the nominal mass and adapts the two boundaries online: at each reset the mass is sampled on a boundary with probability $0.5$ and uniformly in the interior otherwise.
    The mean of the last $20$ returns at each boundary is compared to a running baseline $b$ (mean of the last $50$ episode returns): a boundary is pushed outward by $\Delta=0.1$ kg when performance stays within $20\%$ of $b$, and pulled inward by $\Delta$ when it falls below $50\%$, keeping a minimum gap of $0.05$ kg.
\end{itemize}

The target environment uses a cube mass of 5 kg. However, the randomization range was chosen to be broad enough to include both source and target conditions without being narrowly centered around the target mass, avoiding information leakage from the evaluation environment.

First step of this section was to train the PPO and SAC agents on the source environment with fixed mass, in order to evaluate their behavior and performance on this task.
Both algorithms were implemented using the Stable Baselines3 library \cite{stable-baselines3}, with their hyperparameter configurations selected based on extensive bayesian optimization (taking advantage of \textit{Weigths \& Biases Sweeps \cite{wandb}}).
The final hyperparameters were chosen based on best average reward and success rate across more than 30 training runs for each algorithm, with the best-performing configuration selected for the final evaluation phase.

After training, both PPO and SAC were evaluated on the target environment and compared, aiming to choose the best-performig algorithm as the base on which to apply domain randomization.
The selected algorithm (SAC) was then trained on the source domain under each strategy (none, UDR, ADR) and evaluated deterministically on both the source and target domains over $50$ episodes per seed ($3$ seeds), reporting mean success rate and mean return. The \texttt{none} source$\to$target configuration provides the lower bound and the target$\to$target oracle the upper bound, so UDR and ADR are assessed by how much of that gap they close on the target domain.